% ============================================================================
% PAPER: Sistema Multi-Agente para Generacion Automatica de Preguntas Tipo Test
% Formato: Articulo academico a dos columnas (IEEE style)
% ============================================================================
\documentclass[10pt,twocolumn,letterpaper]{article}

% --- Paquetes ---
\usepackage[utf8]{inputenc}
\usepackage[T1]{fontenc}
\usepackage[spanish]{babel}
\usepackage{times}
\usepackage{lmodern}
\renewcommand{\ttdefault}{lmtt}
\usepackage[margin=1.8cm,top=2cm,bottom=2cm]{geometry}
\usepackage{graphicx}
\usepackage{tikz}
\usetikzlibrary{
  shapes.geometric,
  arrows.meta,
  positioning,
  fit,
  backgrounds,
  calc,
  decorations.pathreplacing,
  patterns,
  shadows,
  chains
}
\usepackage{xcolor}
\usepackage{hyperref}
\usepackage{booktabs}
\usepackage{tabularx}
\usepackage{enumitem}
\usepackage{amsmath}
\usepackage{amssymb}
\usepackage{caption}
\usepackage{float}
\usepackage{fancyhdr}
\usepackage{titlesec}
\usepackage{tcolorbox}
\tcbuselibrary{skins,breakable}
\usepackage{multicol}

% --- Colores personalizados ---
\definecolor{agentZ}{HTML}{4A90D9}
\definecolor{agentB}{HTML}{E8913A}
\definecolor{agentC}{HTML}{50B86C}
\definecolor{agentD}{HTML}{9B59B6}
\definecolor{agentE}{HTML}{E74C3C}
\definecolor{dbcolor}{HTML}{2C3E50}
\definecolor{pipecolor}{HTML}{7F8C8D}
\definecolor{bglight}{HTML}{F8F9FA}
\definecolor{bordercolor}{HTML}{DEE2E6}
\definecolor{codegreen}{HTML}{28A745}
\definecolor{codegray}{HTML}{6C757D}
\definecolor{codeblue}{HTML}{0366D6}
\definecolor{cachegold}{HTML}{F39C12}
\definecolor{reactblue}{HTML}{3498DB}
\definecolor{toolpurple}{HTML}{8E44AD}

% --- Configuracion de secciones ---
\titleformat{\section}{\large\bfseries}{\thesection.}{0.5em}{}
\titleformat{\subsection}{\normalsize\bfseries}{\thesubsection}{0.5em}{}
\titleformat{\subsubsection}{\small\bfseries\itshape}{\thesubsubsection}{0.5em}{}

% --- Header ---
\pagestyle{fancy}
\fancyhf{}
\fancyhead[L]{\small\textit{Sistema Multi-Agente de Generaci\'on de Preguntas}}
\fancyhead[R]{\small\thepage}
\renewcommand{\headrulewidth}{0.4pt}

% --- Code environment sin listings ---
\newtcolorbox{codebox}[1][]{
  colback=bglight, colframe=bordercolor, boxrule=0.5pt, arc=2pt,
  fontupper=\small\ttfamily, left=4pt, right=4pt, top=2pt, bottom=2pt, #1
}

% ============================================================================
\begin{document}

% --- Titulo (one-column) ---
\twocolumn[%
\begin{center}
{\LARGE\bfseries Sistema Multi-Agente para Generaci\'on Autom\'atica
de Preguntas Tipo Test con Control de Calidad}\par
\vspace{10pt}
{\large Documentaci\'on T\'ecnica del Sistema}\par
\vspace{6pt}
{\normalsize Isyfu --- Plataforma de Preparaci\'on de Oposiciones}\par
\vspace{4pt}
{\small Marzo 2026}\par
\vspace{4pt}
\rule{\textwidth}{0.4pt}
\end{center}
\vspace{8pt}
\noindent\textbf{Resumen.}
Se presenta un sistema multi-agente basado en LangGraph para la generaci\'on autom\'atica
de preguntas tipo test a partir de documentos legislativos. El sistema emplea cinco agentes
especializados que operan en un patr\'on Map-Reduce con arquitectura \textbf{StateGraph + ReAct} (Reasoning and Acting, razonamiento y actuaci\'on):
(Z)~fragmentaci\'on coherente con aprendizaje de patrones y an\'alisis de coherencia,
(B)~generaci\'on de preguntas con aprendizaje few-shot (aprendizaje con pocos ejemplos),
8 t\'ecnicas de distracci\'on y
recuperaci\'on multi-chunk basada en FAISS (Facebook AI Similarity Search, b\'usqueda por similitud vectorial),
(C)~evaluaci\'on dual de calidad mediante criterios especializados para oposiciones y m\'etricas Ragas
(Retrieval Augmented Generation Assessment, evaluaci\'on de generaci\'on aumentada por recuperaci\'on),
con m\'ultiples estrategias de validaci\'on (consistencia del tip, validaci\'on de distractores, verificaci\'on de contexto),
(D)~persistencia con deduplicaci\'on sem\'antica mediante embeddings (representaciones vectoriales num\'ericas de texto)
duales (OpenAI/sentence-transformers)
y verificaci\'on post-inserci\'on, y (E)~exportaci\'on a PDF profesional en 3 formatos.
La arquitectura incluye procesamiento paralelo con \texttt{ThreadPoolExecutor}
(pool de hilos para ejecuci\'on concurrente de Python),
m\'ultiples capas de cach\'e (chunks, uso, coordinaci\'on, patrones),
reintentos autom\'aticos con retroalimentaci\'on espec\'ifica, un sistema de configuraci\'on
jer\'arquico basado en Pydantic BaseSettings (librer\'ia de validaci\'on de datos de Python),
y un esquema de base de datos relacional V2
compatible con Supabase (plataforma backend cloud de c\'odigo abierto). El sistema procesa actualmente 23 temas legislativos para
la preparaci\'on de oposiciones a la Guardia Civil espa\~nola.
\vspace{14pt}
]

% ============================================================================
\section{Introducci\'on}
% ============================================================================

La preparaci\'on de oposiciones requiere un volumen masivo de preguntas tipo test de alta
calidad, fundamentadas en legislaci\'on vigente y con distractores plausibles. La creaci\'on
manual de estas preguntas es un proceso costoso, lento y propenso a inconsistencias.

Este documento describe un sistema multi-agente que automatiza el ciclo completo:
desde la ingesta de documentos PDF legislativos hasta la generaci\'on de ex\'amenes en PDF,
pasando por la generaci\'on inteligente de preguntas, validaci\'on de calidad y persistencia
en base de datos.

El sistema se distingue por:
\begin{itemize}[nosep,leftmargin=*]
  \item Arquitectura de 5 agentes con \textbf{StateGraph} (grafo de estados) + \textbf{ReAct} usando LangGraph (framework de agentes basado en grafos de LangChain)
  \item Aprendizaje few-shot con 6 ejemplos reales de ex\'amenes
  \item 8 t\'ecnicas de distracci\'on basadas en an\'alisis de 105 ex\'amenes reales
  \item Validaci\'on dual: criterios especializados + m\'etricas Ragas
  \item 5 estrategias de evaluaci\'on complementarias en Agente~C
  \item Deduplicaci\'on sem\'antica con embeddings duales: OpenAI (API remota) y sentence-transformers (modelos locales de representaci\'on vectorial)
  \item Recuperaci\'on multi-chunk con FAISS para contexto enriquecido
  \item 4 capas de cach\'e para optimizaci\'on de rendimiento
  \item Procesamiento paralelo con \texttt{ThreadPoolExecutor} y cola de escritura
  \item Sistema de configuraci\'on jer\'arquico con Pydantic BaseSettings
  \item Esquema de BD V2 con vistas SQL y compatibilidad Supabase
\end{itemize}

% ============================================================================
\section{Arquitectura General del Sistema}
% ============================================================================

La Figura~\ref{fig:architecture} muestra la arquitectura de alto nivel. El sistema
sigue un patr\'on \textbf{Map-Reduce} (dividir-procesar-combinar): el Agente~Z fragmenta
documentos en chunks (fragmentos de texto) durante la fase ``map'' (divisi\'on),
los Agentes~B y C procesan cada chunk en paralelo (map), y los Agentes~D
y E consolidan resultados durante la fase ``reduce'' (combinaci\'on).

Todos los agentes est\'an implementados como \textbf{StateGraph} (grafo de estados
dirigido) de LangGraph con el patr\'on \textbf{ReAct} (Reasoning + Acting: el modelo
de lenguaje razona, act\'ua ejecutando herramientas, y observa el resultado en un
ciclo iterativo).

\begin{figure*}[t]
\centering
\begin{tikzpicture}[
  node distance=1.2cm and 1.8cm,
  agent/.style={
    rectangle, rounded corners=6pt, minimum width=3.2cm, minimum height=1.4cm,
    text width=3cm, align=center, font=\small\bfseries,
    draw, line width=1pt, drop shadow={shadow xshift=1pt, shadow yshift=-1pt, opacity=0.3}
  },
  data/.style={
    rectangle, rounded corners=3pt, minimum width=2.4cm, minimum height=0.8cm,
    text width=2.2cm, align=center, font=\scriptsize,
    draw, dashed, line width=0.6pt, fill=bglight
  },
  cache/.style={
    rectangle, rounded corners=3pt, minimum width=2cm, minimum height=0.6cm,
    text width=1.8cm, align=center, font=\scriptsize,
    draw, line width=0.6pt, fill=cachegold!15, draw=cachegold
  },
  arrow/.style={-{Stealth[length=6pt]}, line width=1.2pt},
  darrow/.style={-{Stealth[length=5pt]}, line width=0.8pt, dashed},
  label/.style={font=\scriptsize\itshape, text=codegray}
]

% Agentes
\node[agent, fill=agentZ!15, draw=agentZ] (Z) {Agente Z\\{\scriptsize Fragmentaci\'on\\ y Reescritura}};
\node[agent, fill=agentB!15, draw=agentB, right=2.5cm of Z] (B) {Agente B\\{\scriptsize Generaci\'on\\de Preguntas}};
\node[agent, fill=agentC!15, draw=agentC, right=2.5cm of B] (C) {Agente C\\{\scriptsize Evaluaci\'on\\de Calidad}};
\node[agent, fill=agentD!15, draw=agentD, below=2cm of B] (D) {Agente D\\{\scriptsize Persistencia\\y Deduplicaci\'on}};
\node[agent, fill=agentE!15, draw=agentE, right=2.5cm of D] (E) {Agente E\\{\scriptsize Generaci\'on\\de PDFs}};

% Datos entrada/salida
\node[data, left=1.5cm of Z] (input) {PDF / TXT / MD\\{\tiny input\_docs/}};
\node[data, below=0.6cm of E] (pdf) {Ex\'amenes PDF\\{\tiny output/pdfs/}};

% Caches
\node[cache, below=0.6cm of Z] (cache1) {JSON Cache\\{\tiny rewritten/}};
\node[cache, below=0.3cm of cache1] (cache2) {Pattern Registry\\{\tiny patterns.json}};
\node[cache, below=0.6cm of D] (db) {SQLite DB V2\\{\tiny questions.db}};

% FAISS
\node[rectangle, rounded corners=3pt, draw=reactblue, fill=reactblue!10,
      minimum width=2cm, minimum height=0.6cm,
      font=\scriptsize\bfseries, align=center, below=0.3cm of B] (faiss)
  {FAISS Index\\{\tiny embeddings}};

% Flujo principal
\draw[arrow, color=agentZ] (input) -- (Z);
\draw[arrow, color=agentZ] (Z) -- node[above, label] {chunks} (B);
\draw[arrow, color=agentB] (B) -- node[above, label] {preguntas} (C);
\draw[arrow, color=agentC] (C) -- node[right, label, pos=0.3] {validadas} (D);
\draw[arrow, color=agentD] (D) -- node[above, label] {IDs} (E);

% Flujo secundario
\draw[darrow, color=agentZ] (Z) -- (cache1);
\draw[darrow, color=agentZ] (Z.south) -- ++(0,-0.3) -| (cache2);
\draw[darrow, color=agentD] (D) -- (db);
\draw[darrow, color=agentE] (E) -- (pdf);
\draw[darrow, color=reactblue] (B) -- (faiss);

% Retry loop
\draw[arrow, color=agentE!70!black, line width=0.9pt]
  (C.south) -- ++(0,-0.5) -| node[below, label, pos=0.25] {retry con feedback (m\'ax 3)} (B.south);

% Background pipeline box
\begin{scope}[on background layer]
  \node[fit=(Z)(B)(C)(D)(E), inner sep=12pt, rounded corners=10pt,
        draw=pipecolor, line width=0.6pt, dashed, fill=white] {};
\end{scope}

\end{tikzpicture}
\caption{Arquitectura general del sistema multi-agente. Las flechas s\'olidas representan
el flujo principal del pipeline; las discontinuas, almacenamiento intermedio.
Se muestran las capas de cach\'e (dorado), el \'indice FAISS (azul) y el bucle de
retry entre C y B con retroalimentaci\'on espec\'ifica.}
\label{fig:architecture}
\end{figure*}

% ============================================================================
\section{Patr\'on ReAct y StateGraph}
\label{sec:react}
% ============================================================================

Cada agente del sistema est\'a implementado como un \textbf{StateGraph} de LangGraph
que sigue el patr\'on \textbf{ReAct} (Reasoning + Acting). Este patr\'on permite al
LLM razonar sobre el estado actual, decidir qu\'e herramienta ejecutar, observar el
resultado, y repetir hasta completar la tarea.

La Figura~\ref{fig:react} ilustra el ciclo ReAct utilizado en todos los agentes.

\begin{figure}[H]
\centering
\begin{tikzpicture}[
  node distance=1cm,
  state/.style={
    circle, draw, line width=1pt, minimum size=1.4cm,
    font=\scriptsize\bfseries, align=center, inner sep=2pt
  },
  arrow/.style={-{Stealth[length=5pt]}, line width=1pt}
]

% States
\node[state, fill=reactblue!15, draw=reactblue] (think) at (0,0) {Reason\\{\tiny pensar}};
\node[state, fill=toolpurple!15, draw=toolpurple] (act) at (3,0) {Act\\{\tiny ejecutar}};
\node[state, fill=agentC!15, draw=agentC] (observe) at (1.5,-2) {Observe\\{\tiny observar}};

% Decision
\node[diamond, draw=agentE, fill=agentE!10, minimum size=0.8cm,
      font=\tiny\bfseries, inner sep=1pt] (done) at (5.5,0) {Fin?};
\node[rectangle, rounded corners=3pt, draw=codegreen, fill=codegreen!15,
      font=\scriptsize\bfseries, minimum width=1.2cm] (output) at (5.5,-2) {Output};

% Arrows
\draw[arrow, color=reactblue] (think) -- node[above, font=\tiny] {tool\_call} (act);
\draw[arrow, color=toolpurple] (act) -- node[right, font=\tiny] {resultado} (observe);
\draw[arrow, color=agentC] (observe) -- node[left, font=\tiny] {actualizar\\estado} (think);
\draw[arrow, color=agentE] (act) -- (done);
\draw[arrow, color=codegreen] (done) -- node[right, font=\tiny] {s\'i} (output);
\draw[arrow, color=pipecolor] (done.north) -- ++(0,0.6) -| node[above, font=\tiny, pos=0.25] {no} (think.north);

% Input
\node[rectangle, rounded corners=3pt, draw=pipecolor, fill=bglight,
      font=\scriptsize, minimum width=1.4cm] (input) at (-2,0) {GraphState};
\draw[arrow, color=pipecolor] (input) -- (think);

\end{tikzpicture}
\caption{Patr\'on ReAct implementado en cada agente. El LLM razona (Reason),
ejecuta una herramienta (Act), observa el resultado (Observe), y decide si
continuar o finalizar. Implementado con \texttt{langgraph.prebuilt.create\_react\_agent}.}
\label{fig:react}
\end{figure}

Cada agente define sus propias herramientas (\textit{tools}) que el LLM puede invocar:

\begin{itemize}[nosep,leftmargin=*]
  \item \textbf{Agente Z}: \texttt{load\_document}, \texttt{filter\_metadata},
        \texttt{analyze\_structure}, \texttt{create\_chunks}, \texttt{evaluate\_coherence}
  \item \textbf{Agente B}: \texttt{generate\_questions}, \texttt{retrieve\_context},
        \texttt{evaluate\_sufficiency}
  \item \textbf{Agente C}: \texttt{evaluate\_specialized}, \texttt{evaluate\_ragas},
        \texttt{check\_tip\_consistency}, \texttt{validate\_distractors}
  \item \textbf{Agente D}: \texttt{compute\_embedding}, \texttt{check\_duplicate},
        \texttt{insert\_question}, \texttt{verify\_insertion}
  \item \textbf{Agente E}: \texttt{query\_questions}, \texttt{generate\_pdf},
        \texttt{create\_answer\_sheet}
\end{itemize}

La \textbf{inyecci\'on de contexto en tiempo de ejecuci\'on} (runtime context injection)
permite a las herramientas del StateGraph capturar
y modificar el estado compartido (\texttt{GraphState}) durante la ejecuci\'on, lo que
habilita la comunicaci\'on impl\'icita entre pasos del grafo.

% ============================================================================
\section{Agentes del Sistema}
% ============================================================================

\subsection{Agente Z: Fragmentaci\'on Coherente}

El Agente~Z es responsable de transformar documentos legislativos en bruto en
\textbf{chunks sem\'anticamente coherentes}. Es el agente m\'as complejo del sistema
(2584 LOC) e incluye m\'ultiples subsistemas especializados.

\subsubsection{Pipeline de Procesamiento}

\begin{enumerate}[nosep,leftmargin=*]
  \item \textbf{Carga del documento}: Soporta PDF (pypdf, librer\'ia de lectura de PDF),
        TXT y MD (Markdown) con extracci\'on autom\'atica de metadatos (tipo, tama\~no, p\'aginas)
  \item \textbf{Filtrado de metadatos}: El \texttt{PDFMetadataFilter} elimina ruido
        con detecci\'on de emails, URLs, tel\'efonos, ISBN, dep\'osito legal,
        headers/footers y branding
  \item \textbf{An\'alisis estructural}: Identifica art\'iculos, cap\'itulos, secciones
        y t\'itulos del documento legal
  \item \textbf{Coordinaci\'on de chunks}: An\'alisis mediante LLM (Large Language Model,
        modelo de lenguaje de gran tama\~no) o heur\'istico de l\'imites
        de fragmentaci\'on (\texttt{coordination\_mode: llm|heuristic})
  \item \textbf{Fragmentaci\'on}: Genera chunks con tama\~no configurable (por defecto
        1000 caracteres, 200 de solapamiento)
  \item \textbf{An\'alisis de coherencia}: Clasifica cada chunk y decide acci\'on
  \item \textbf{Evaluaci\'on de sustantividad}: Ratio contenido real vs.\ ruido
        (\texttt{min\_substantive\_ratio: 0.3})
  \item \textbf{Guardado incremental}: Cada $N$ chunks para prevenir p\'erdida
\end{enumerate}

\subsubsection{An\'alisis de Coherencia de Chunks}

La Figura~\ref{fig:coherence} muestra el sistema de clasificaci\'on de coherencia.

\begin{figure}[H]
\centering
\begin{tikzpicture}[
  type/.style={
    rectangle, rounded corners=3pt, draw, line width=0.6pt,
    minimum width=2.8cm, minimum height=0.55cm,
    text width=2.6cm, align=center, font=\scriptsize
  },
  action/.style={
    rectangle, rounded corners=3pt, draw, line width=0.6pt,
    minimum width=2cm, minimum height=0.55cm,
    text width=1.8cm, align=center, font=\scriptsize\bfseries
  },
  arrow/.style={-{Stealth[length=4pt]}, line width=0.6pt}
]

% Chunk Types
\node[font=\scriptsize\bfseries, color=agentZ] at (-1.6,1.8) {ChunkType};
\node[type, fill=agentZ!20] (t1) at (-1.6,1.2) {COMPLETE\_ARTICLE};
\node[type, fill=agentZ!15] (t2) at (-1.6,0.5) {COMPLETE\_SECTION};
\node[type, fill=agentZ!10] (t3) at (-1.6,-0.2) {PARTIAL\_CONTENT};
\node[type, fill=agentZ!8] (t4) at (-1.6,-0.9) {METADATA};
\node[type, fill=agentE!10] (t5) at (-1.6,-1.6) {GARBAGE};

% Chunk Actions
\node[font=\scriptsize\bfseries, color=agentC] at (2,1.8) {ChunkAction};
\node[action, fill=codegreen!15, draw=codegreen] (a1) at (2,1.2) {KEEP};
\node[action, fill=agentB!15, draw=agentB] (a2) at (2,0.5) {MERGE\_NEXT};
\node[action, fill=agentB!10, draw=agentB] (a3) at (2,-0.2) {MERGE\_PREV};
\node[action, fill=agentE!15, draw=agentE] (a4) at (2,-1.25) {DISCARD};

% Arrows
\draw[arrow, color=codegreen] (t1.east) -- (a1.west);
\draw[arrow, color=codegreen] (t2.east) -- (a1.west);
\draw[arrow, color=agentB] (t3.east) -- (a2.west);
\draw[arrow, color=agentB] (t3.east) -- (a3.west);
\draw[arrow, color=agentE] (t4.east) -- (a4.west);
\draw[arrow, color=agentE] (t5.east) -- (a4.west);

% Confidence
\node[font=\tiny\itshape, color=codegray, text width=2.8cm, align=center] at (0.2,-2.3)
  {Cada an\'alisis incluye\\confidence score $[0,1]$};

\end{tikzpicture}
\caption{Sistema de an\'alisis de coherencia de chunks. Cada chunk se clasifica
por tipo y se le asigna una acci\'on. Los chunks de tipo \texttt{PARTIAL\_CONTENT}
se fusionan con adyacentes; \texttt{METADATA} y \texttt{GARBAGE} se descartan.}
\label{fig:coherence}
\end{figure}

\subsubsection{Registro de Patrones Aprendidos}

El \texttt{RewritePatternRegistry} implementa un sistema de \textbf{aprendizaje
adaptativo} que registra patrones regex (expresiones regulares, reglas de b\'usqueda
de texto) descubiertos durante la limpieza de documentos:

\begin{itemize}[nosep,leftmargin=*]
  \item Aprende autom\'aticamente patrones de ruido recurrente
  \item Persiste patrones en \texttt{patterns.json} entre ejecuciones
  \item Filtra por \texttt{min\_hits}, \texttt{min\_len}, \texttt{max\_len}
  \item Reutiliza patrones aprendidos en documentos nuevos
\end{itemize}

\begin{tcolorbox}[colback=agentZ!5, colframe=agentZ, title={\small Modelo de Chunk},
  fonttitle=\bfseries\small, boxrule=0.5pt, arc=3pt]
{\small\ttfamily
class Chunk(BaseModel):\\
\hspace*{1em}chunk\_id: str\\
\hspace*{1em}content: str\\
\hspace*{1em}clean\_content: Optional[str]\\
\hspace*{1em}source\_document: str\\
\hspace*{1em}page: int\\
\hspace*{1em}token\_count: int\\
\hspace*{1em}is\_substantive: bool\\
\hspace*{1em}substantive\_ratio: float\\
\hspace*{1em}embedding\_indexed: bool\\
\hspace*{1em}related\_chunk\_ids: List[str]\\[4pt]
\hspace*{1em}def get\_content\_for\_generation()\\
\hspace*{1em}def get\_effective\_token\_count()
}
\end{tcolorbox}

% ---
\subsection{Agente B: Generaci\'on de Preguntas}

El Agente~B constituye el n\'ucleo generativo del sistema (2143 LOC). Implementado como
\texttt{StateGraph} con aprendizaje few-shot, genera preguntas tipo test
de calidad profesional.

\subsubsection{Aprendizaje Few-Shot}
Se incorporan \textbf{6 ejemplos reales} de ex\'amenes oficiales de la Guardia Civil
en formato JSON (JavaScript Object Notation, formato est\'andar de intercambio de datos)
estructurado (\texttt{agent\_b\_generation.json}) directamente en el prompt (instrucci\'on
al modelo de lenguaje) del sistema. Estos ejemplos establecen el formato, nivel de dificultad y
estilo esperado.

\subsubsection{T\'ecnicas de Distracci\'on}
Bas\'andose en el an\'alisis de 105 ex\'amenes reales, se han documentado 8 t\'ecnicas
para crear distractores plausibles (Tabla~\ref{tab:distractors}).

\begin{table}[H]
\centering
\caption{T\'ecnicas de distracci\'on documentadas}
\label{tab:distractors}
\scriptsize
\begin{tabularx}{\columnwidth}{cX}
\toprule
\textbf{\#} & \textbf{T\'ecnica} \\
\midrule
1 & Cambio de conjunci\'on/disyunci\'on (Y $\leftrightarrow$ O) \\
2 & Manipulaci\'on de plazos temporales \\
3 & Eliminaci\'on de excepciones legales \\
4 & Atribuci\'on incorrecta de funciones \\
5 & Confusi\'on de nivel normativo (ley vs.\ reglamento) \\
6 & Falsa limitaci\'on de \'ambito \\
7 & Alteraci\'on de cifras o cantidades \\
8 & Contexto correcto con final modificado \\
\bottomrule
\end{tabularx}
\end{table}

\subsubsection{Estrategias de Expansi\'on de Contexto}

El Agente~B implementa dos estrategias complementarias para enriquecer el contexto
de generaci\'on (Figura~\ref{fig:expansion}):

\begin{figure}[H]
\centering
\begin{tikzpicture}[
  chunk/.style={
    rectangle, rounded corners=2pt, draw, line width=0.5pt,
    minimum width=1cm, minimum height=0.6cm,
    font=\tiny, align=center
  },
  strategy/.style={
    rectangle, rounded corners=4pt, draw, line width=0.8pt,
    minimum width=3.6cm, minimum height=0.7cm,
    font=\scriptsize\bfseries, align=center
  },
  arr/.style={-{Stealth[length=4pt]}, line width=0.6pt}
]

% Strategy 1: Sequential
\node[strategy, fill=agentB!10, draw=agentB] (s1) at (0,0) {Expansi\'on Secuencial};
\node[chunk, fill=pipecolor!10] (p1) at (-1.5,-0.9) {Chunk\\N-1};
\node[chunk, fill=agentB!25] (p2) at (0,-0.9) {Chunk\\N};
\node[chunk, fill=pipecolor!10] (p3) at (1.5,-0.9) {Chunk\\N+1};

\draw[arr, color=agentB] (s1) -- (p2);
\draw[arr, color=pipecolor, dashed] (p2) -- (p1);
\draw[arr, color=pipecolor, dashed] (p2) -- (p3);

% Strategy 2: Semantic
\node[strategy, fill=reactblue!10, draw=reactblue] (s2) at (0,-2.4) {Expansi\'on Sem\'antica};
\node[chunk, fill=reactblue!25] (f1) at (-1.5,-3.3) {Chunk\\Art.5};
\node[chunk, fill=reactblue!40] (f2) at (0,-3.3) {Chunk\\Art.7};
\node[chunk, fill=reactblue!20] (f3) at (1.5,-3.3) {Chunk\\Art.12};

\node[rectangle, rounded corners=2pt, draw=reactblue, fill=reactblue!8,
      font=\tiny\bfseries, minimum width=1.4cm] (faiss) at (0,-4.1) {FAISS};

\draw[arr, color=reactblue] (s2) -- (f2);
\draw[arr, color=reactblue, dashed] (faiss) -- (f1);
\draw[arr, color=reactblue, dashed] (faiss) -- (f3);
\draw[arr, color=reactblue] (f2) -- (faiss);

% Labels
\node[font=\tiny\itshape, color=codegray] at (3.2,-0.9) {merge prev/next};
\node[font=\tiny\itshape, color=codegray] at (3.2,-3.3) {top-k similares};

\end{tikzpicture}
\caption{Dos estrategias de expansi\'on de contexto. La \textbf{secuencial} fusiona
chunks adyacentes; la \textbf{sem\'antica} recupera chunks similares v\'ia FAISS
independientemente de su posici\'on en el documento.}
\label{fig:expansion}
\end{figure}

\subsubsection{Modelo Interno de Pregunta}

Cada pregunta generada incluye metadatos de trazabilidad:

\begin{tcolorbox}[colback=agentB!5, colframe=agentB, boxrule=0.5pt, arc=3pt,
  title={\small QuestionData}, fonttitle=\bfseries\small]
{\scriptsize\ttfamily
source\_chunk\_ids: List[str]\\
source\_chunk\_indices: List[int]\\
source\_doc\_ids: List[str]\\
context\_strategy: str ~~\# single|multi\\
generation\_time: float\\
reasoning\_summary: str\\
distractor\_techniques: List[int]
}
\end{tcolorbox}

% ---
\subsection{Agente C: Evaluaci\'on de Calidad}

El Agente~C implementa un \textbf{sistema dual de validaci\'on} (1751 LOC) que combina
m\'ultiples estrategias de evaluaci\'on. Es el agente con m\'as prompts especializados
del sistema.

\subsubsection{Estrategias de Evaluaci\'on}

La Figura~\ref{fig:evaluation} muestra las 5 estrategias de evaluaci\'on disponibles.

\begin{figure}[H]
\centering
\begin{tikzpicture}[
  eval/.style={
    rectangle, rounded corners=3pt, draw, line width=0.6pt,
    minimum width=3.4cm, minimum height=0.55cm,
    text width=3.2cm, align=center, font=\scriptsize
  },
  prio/.style={
    rectangle, rounded corners=2pt, draw, line width=0.5pt,
    minimum width=1.2cm, minimum height=0.4cm,
    font=\tiny\bfseries, align=center
  },
  arr/.style={-{Stealth[length=4pt]}, line width=0.6pt}
]

% Central node
\node[rectangle, rounded corners=5pt, draw=agentC, fill=agentC!15,
      line width=1pt, font=\small\bfseries, minimum width=2.5cm,
      minimum height=0.7cm] (center) at (0,0.5) {Agente C};

% Evaluations
\node[eval, fill=agentC!10] (e1) at (0,-0.6) {Eval.\ Especializada Oposiciones};
\node[eval, fill=agentC!8] (e2) at (0,-1.4) {M\'etricas Ragas (F + R)};
\node[eval, fill=agentC!6] (e3) at (0,-2.2) {Consistencia del Tip};
\node[eval, fill=agentC!5] (e4) at (0,-3.0) {Validaci\'on de Distractores};
\node[eval, fill=agentC!4] (e5) at (0,-3.8) {Verificaci\'on de Contexto};

% Priority labels
\node[prio, fill=agentE!20, draw=agentE] at (3.3,-0.6) {PRINCIPAL};
\node[prio, fill=agentB!15, draw=agentB] at (3.3,-1.4) {COMPLEM.};
\node[prio, fill=pipecolor!15, draw=pipecolor] at (3.3,-2.2) {OPCIONAL};
\node[prio, fill=pipecolor!15, draw=pipecolor] at (3.3,-3.0) {OPCIONAL};
\node[prio, fill=pipecolor!15, draw=pipecolor] at (3.3,-3.8) {OPCIONAL};

% Arrows
\foreach \i in {1,...,5} {
  \draw[arr, color=agentC] (center.south) -- (e\i.north);
}

% Fast mode note
\node[rectangle, rounded corners=2pt, draw=cachegold, fill=cachegold!10,
      font=\tiny\itshape, text width=3cm, align=center] at (0,-4.6)
  {Fast mode: solo eval.\ especializada\\+ Ragas (skip opcionales)};

\end{tikzpicture}
\caption{Cinco estrategias de evaluaci\'on del Agente~C. Cada una tiene su propio
prompt especializado. En modo r\'apido, solo se ejecutan las dos principales.}
\label{fig:evaluation}
\end{figure}

\subsubsection{Evaluaci\'on Especializada (Principal)}
Cinco criterios ponderados por criticidad:

\begin{itemize}[nosep,leftmargin=*]
  \item[\textbf{C}] \textbf{Referencias legales expl\'icitas} --- La pregunta debe
        citar art\'iculos o leyes espec\'ificas
  \item[\textbf{C}] \textbf{Precisi\'on t\'ecnica absoluta} --- Sin ambig\"uedades ni
        errores jur\'idicos
  \item[\textbf{A}] \textbf{Distractores plausibles} --- Las opciones incorrectas
        deben ser cre\'ibles
  \item[\textbf{A}] \textbf{Retroalimentaci\'on estructurada} --- El ``tip'' debe
        explicar por qu\'e la respuesta es correcta (40--200 palabras)
  \item[\textbf{M}] \textbf{Formato oficial} --- 4 opciones, sin ambig\"uedad en
        la correcta
\end{itemize}

{\scriptsize (C = Cr\'itico, A = Alto, M = Medio)}

\subsubsection{M\'etricas Ragas (Complementaria)}
\begin{itemize}[nosep,leftmargin=*]
  \item \textbf{Faithfulness} (fidelidad): \textquestiondown Es la respuesta fiel al contexto original? ($0$--$1$)
  \item \textbf{Answer Relevancy} (relevancia de la respuesta): \textquestiondown Es la respuesta relevante a la pregunta? ($0$--$1$)
\end{itemize}

\subsubsection{Evaluaciones Opcionales}

\begin{itemize}[nosep,leftmargin=*]
  \item \textbf{Tip Consistency} (consistencia de la explicaci\'on): Verifica que el tip
        (explicaci\'on did\'actica de la respuesta correcta) sea coherente con la pregunta
        y la respuesta correcta (\texttt{agent\_c\_tip\_consistency\_system.txt})
  \item \textbf{Tip Supports Answer}: Valida que el tip apoye expl\'icitamente la
        respuesta marcada como correcta (\texttt{agent\_c\_tip\_supports\_answer\_system.txt})
  \item \textbf{Distractor Validation}: Comprueba que los distractores usen las
        t\'ecnicas documentadas (\texttt{agent\_c\_validate\_distractor\_system.txt})
  \item \textbf{Context Check} (verificaci\'on de contexto): Eval\'ua si el contexto
        (fragmento de texto origen) es suficiente para la pregunta generada,
        con detecci\'on de nivel de complejidad
\end{itemize}

\subsubsection{Clasificaci\'on de Resultados}

La Figura~\ref{fig:quality} ilustra el sistema de clasificaci\'on tripartito.

\begin{figure}[H]
\centering
\begin{tikzpicture}[
  box/.style={rectangle, rounded corners=4pt, minimum width=2.8cm, minimum height=0.9cm,
    text width=2.6cm, align=center, font=\scriptsize\bfseries, draw, line width=0.8pt},
  score/.style={font=\tiny, text=codegray}
]

% Boxes
\node[box, fill=green!15, draw=green!60!black] (pass) at (0,0) {AUTO\_PASS\\{\tiny Aprobado}};
\node[box, fill=yellow!15, draw=yellow!60!black] (review) at (0,-1.3) {MANUAL\_REVIEW\\{\tiny Revisi\'on humana}};
\node[box, fill=red!15, draw=red!60!black] (fail) at (0,-2.6) {AUTO\_FAIL\\{\tiny Regenerar}};

% Scores
\node[score, right=0.3cm of pass] {$F \geq 0.85$ \textbf{y} $R \geq 0.85$};
\node[score, right=0.3cm of review] {$0.60 \leq F < 0.85$};
\node[score, right=0.3cm of fail] {$F < 0.60$ \textbf{\'o} $R < 0.60$};

% Brace
\draw[decorate, decoration={brace, amplitude=5pt, mirror}, line width=0.6pt]
  (-1.6,0.5) -- (-1.6,-3.1) node[midway, left=8pt, font=\scriptsize, align=right]
  {Umbral\\configurable};

\end{tikzpicture}
\caption{Sistema de clasificaci\'on de calidad. $F$ = Faithfulness, $R$ = Answer Relevancy.
Los umbrales son configurables por perfil.}
\label{fig:quality}
\end{figure}

Las preguntas clasificadas como \texttt{auto\_fail} se reenv\'ian al Agente~B con
retroalimentaci\'on espec\'ifica (motivo del fallo, criterios incumplidos),
permitiendo hasta \textbf{3 reintentos} por chunk. Adicionalmente, existe un
l\'imite de \textbf{2 reintentos por pregunta individual} (\texttt{MAX\_RETRIES\_PER\_QUESTION}).

% ---
\subsection{Agente D: Persistencia y Deduplicaci\'on}

El Agente~D gestiona la inserci\'on en base de datos con deduplicaci\'on sem\'antica
(detecci\'on de preguntas duplicadas por significado, no solo por texto id\'entico)
y verificaci\'on post-inserci\'on:

\begin{enumerate}[nosep,leftmargin=*]
  \item Genera embedding (vector num\'erico que representa el significado) de
        \texttt{question + correct\_answer}
  \item Calcula similitud coseno (medida matem\'atica de cercan\'ia entre vectores, $0$--$1$)
        contra preguntas existentes
  \item Si similitud $> 0.85$: marca como duplicado (\texttt{is\_duplicate = true})
  \item Si es \'unica: inserta en SQLite con todos los metadatos
  \item \textbf{Verificaci\'on}: \texttt{\_verify\_inserted()} comprueba integridad
  \item Registra batch con estad\'isticas (total, \'unicas, duplicadas, scores medios)
\end{enumerate}

\subsubsection{Sistema de Embeddings Dual}

La Figura~\ref{fig:embeddings} muestra el sistema dual de embeddings.

\begin{figure}[H]
\centering
\begin{tikzpicture}[
  provider/.style={
    rectangle, rounded corners=4pt, draw, line width=0.8pt,
    minimum width=3cm, minimum height=1.2cm,
    text width=2.8cm, align=center, font=\scriptsize
  },
  arr/.style={-{Stealth[length=4pt]}, line width=0.7pt}
]

% Central
\node[rectangle, rounded corners=5pt, draw=agentD, fill=agentD!10,
      line width=1pt, font=\small\bfseries, minimum width=3.2cm,
      minimum height=0.7cm] (service) at (0,0) {EmbeddingService};

% Providers
\node[provider, fill=reactblue!10, draw=reactblue] (openai) at (-2.2,-1.8)
  {\textbf{OpenAI}\\{\tiny text-embedding-3-small}\\{\tiny API remota}};
\node[provider, fill=codegreen!10, draw=codegreen] (st) at (2.2,-1.8)
  {\textbf{sentence-transformers}\\{\tiny all-MiniLM-L6-v2}\\{\tiny Local / GPU}};

% FAISS
\node[rectangle, rounded corners=3pt, draw=cachegold, fill=cachegold!10,
      font=\scriptsize\bfseries, minimum width=2.4cm, minimum height=0.6cm]
  (faiss) at (0,-3.4) {FAISS Index};

% Arrows
\draw[arr, color=reactblue] (service) -- (openai);
\draw[arr, color=codegreen] (service) -- (st);
\draw[arr, color=cachegold] (openai.south) -- ++(0,-0.4) -| (faiss);
\draw[arr, color=cachegold] (st.south) -- ++(0,-0.4) -| (faiss);

% Config
\node[font=\tiny\ttfamily, color=codegray, text width=3cm, align=center] at (0,-4.2)
  {EMBEDDING\_PROVIDER=openai|st\\SIMILARITY\_THRESHOLD=0.7};

\end{tikzpicture}
\caption{Sistema dual de embeddings. Configurable entre OpenAI (remoto, mayor calidad)
y sentence-transformers (local, sin coste). Ambos alimentan FAISS para b\'usqueda
eficiente.}
\label{fig:embeddings}
\end{figure}

% ---
\subsection{Agente E: Generaci\'on de PDFs}

Implementado con ReportLab (librer\'ia Python para creaci\'on de documentos PDF)
con 1159 LOC (l\'ineas de c\'odigo), genera tres formatos de salida:

\begin{itemize}[nosep,leftmargin=*]
  \item \textbf{EXAM}: Solo preguntas con opciones, sin soluciones
  \item \textbf{WITH\_SOLUTIONS}: Incluye tips y explicaciones detalladas
  \item \textbf{ANSWERS\_ONLY}: Plantilla de correcci\'on r\'apida
\end{itemize}

Caracter\'isticas avanzadas:
\begin{itemize}[nosep,leftmargin=*]
  \item Normalizaci\'on de prefijos de opciones (A/B/C/D)
  \item Normalizaci\'on de tips: reemplaza referencias ``opci\'on A'' por texto real
  \item Truncado inteligente de texto (\texttt{max\_len=600})
  \item Histogramas de dificultad con Matplotlib
  \item Agrupaci\'on por tema con estad\'isticas
  \item Tablas estilizadas con ReportLab
\end{itemize}

% ============================================================================
\section{Recuperaci\'on Multi-Chunk con FAISS}
% ============================================================================

Una innovaci\'on clave del sistema es la \textbf{recuperaci\'on multi-chunk} para
generar preguntas que requieren contexto de m\'ultiples art\'iculos o secciones.
El \texttt{ChunkRetrieverService} orquesta este proceso.

\begin{figure}[H]
\centering
\begin{tikzpicture}[
  chunk/.style={
    rectangle, rounded corners=3pt, draw, line width=0.6pt,
    minimum width=1.2cm, minimum height=0.9cm,
    font=\tiny, align=center
  },
  enriched/.style={
    rectangle, rounded corners=5pt, draw=agentB, line width=1pt,
    minimum width=3.4cm, minimum height=1.4cm,
    font=\scriptsize, align=center, fill=agentB!5
  }
]

% Chunks
\node[chunk, fill=agentZ!15] (c1) at (0,0) {Chunk 1\\Art. 5};
\node[chunk, fill=agentZ!15] (c2) at (1.4,0) {Chunk 2\\Art. 6};
\node[chunk, fill=agentZ!30] (c3) at (2.8,0) {Chunk 3\\Art. 7\\(primario)};
\node[chunk, fill=agentZ!15] (c4) at (4.2,0) {Chunk 4\\Art. 8};
\node[chunk, fill=agentZ!15] (c5) at (5.6,0) {Chunk 5\\Art. 9};

% FAISS
\node[rectangle, rounded corners=3pt, draw=pipecolor, fill=pipecolor!10,
      font=\tiny\bfseries, minimum width=2cm] (faiss) at (2.8,-1.5) {FAISS Index};

% Enriched
\node[enriched] (ec) at (2.8,-3.3) {\textbf{EnrichedContext}\\[2pt]
  Primario: Art. 7\\Relacionados: Art. 6, 8\\Total tokens: 1847};

% Sufficiency
\node[rectangle, rounded corners=3pt, draw=agentC, fill=agentC!10,
      font=\tiny\bfseries, minimum width=2.2cm] (suf) at (2.8,-4.6)
  {ContextSufficiency\\{\tiny sufficient: true}};

% Arrows
\foreach \c in {c1,c2,c3,c4,c5} {
  \draw[-{Stealth[length=3pt]}, line width=0.4pt, color=pipecolor] (\c) -- (faiss);
}
\draw[-{Stealth[length=4pt]}, line width=0.8pt, color=agentB] (faiss) -- (ec);
\draw[-{Stealth[length=4pt]}, line width=0.6pt, color=agentC] (ec) -- (suf);

% Similarity annotations
\draw[{Stealth[length=3pt]}-, line width=0.5pt, color=agentB!60]
  (c2.south east) -- ++(0.2,-0.3) node[right, font=\tiny] {0.91};
\draw[{Stealth[length=3pt]}-, line width=0.5pt, color=agentB!60]
  (c4.south west) -- ++(-0.2,-0.3) node[left, font=\tiny] {0.88};

\end{tikzpicture}
\caption{Sistema de recuperaci\'on multi-chunk con FAISS. El chunk primario se
enriquece con los $k$ chunks m\'as similares sem\'anticamente. Se eval\'ua la
suficiencia del contexto antes de generar.}
\label{fig:multichunk}
\end{figure}

El flujo completo del \texttt{ChunkRetrieverService}:

\begin{enumerate}[nosep,leftmargin=*]
  \item Filtra chunks no sustantivos ($ratio < 0.3$)
  \item Genera embeddings con el proveedor configurado (OpenAI o sentence-transformers)
  \item Indexa en FAISS para b\'usqueda vectorial eficiente
  \item Para cada chunk primario, recupera m\'aximo 3 relacionados
        (\texttt{max\_related\_chunks: 3})
  \item Eval\'ua \texttt{ContextSufficiency}: suficiente, nivel de complejidad,
        chunks adicionales necesarios
  \item Retorna \texttt{EnrichedContext}: chunk primario + relacionados +
        contenido combinado + total tokens + temas cubiertos
  \item Registra \texttt{context\_strategy} (\texttt{single} o \texttt{multi})
        en metadatos de la pregunta
\end{enumerate}

% ============================================================================
\section{Pipeline de Datos}
% ============================================================================

La Figura~\ref{fig:pipeline} muestra el flujo de datos detallado a trav\'es del
sistema, incluyendo las transformaciones en cada etapa.

\begin{figure*}[t]
\centering
\begin{tikzpicture}[
  phase/.style={
    rectangle, rounded corners=5pt, minimum width=2.6cm, minimum height=2.4cm,
    text width=2.4cm, align=center, font=\scriptsize, draw, line width=0.8pt
  },
  io/.style={
    rectangle, rounded corners=3pt,
    minimum width=1.8cm, minimum height=0.6cm,
    text width=1.6cm, align=center, font=\tiny, draw, line width=0.6pt, fill=bglight
  },
  store/.style={
    rectangle, rounded corners=3pt, draw, minimum height=0.8cm, minimum width=1.4cm,
    font=\tiny, align=center, line width=0.6pt, fill=bglight
  },
  cachenode/.style={
    rectangle, rounded corners=3pt, draw=cachegold, minimum height=0.6cm,
    minimum width=1.4cm, font=\tiny, align=center, line width=0.6pt,
    fill=cachegold!10
  },
  arr/.style={-{Stealth[length=5pt]}, line width=1pt},
  darr/.style={-{Stealth[length=4pt]}, line width=0.6pt, dashed}
]

% Fases
\node[phase, fill=agentZ!10, draw=agentZ] (p1) at (0,0)
  {\textbf{FASE 1}\\[2pt] Ingesta y\\Fragmentaci\'on\\[4pt]
   {\tiny Agente Z}\\[2pt]
   {\tiny Carga PDF}\\
   {\tiny Filtrado metadata}\\
   {\tiny An\'alisis coherencia}\\
   {\tiny Chunks coherentes}};

\node[phase, fill=agentB!10, draw=agentB] (p2) at (3.5,0)
  {\textbf{FASE 2}\\[2pt] Generaci\'on\\Inteligente\\[4pt]
   {\tiny Agente B}\\[2pt]
   {\tiny Few-shot learning}\\
   {\tiny Multi-chunk FAISS}\\
   {\tiny 8 t\'ecn.\ distracci\'on}\\
   {\tiny 2 estrat.\ expansi\'on}};

\node[phase, fill=agentC!10, draw=agentC] (p3) at (7,0)
  {\textbf{FASE 3}\\[2pt] Validaci\'on\\de Calidad\\[4pt]
   {\tiny Agente C}\\[2pt]
   {\tiny 5 eval.\ especializadas}\\
   {\tiny M\'etricas Ragas}\\
   {\tiny Clasificaci\'on 3-v\'ias}\\
   {\tiny Fast mode opcional}};

\node[phase, fill=agentD!10, draw=agentD] (p4) at (10.5,0)
  {\textbf{FASE 4}\\[2pt] Persistencia\\y Dedup\\[4pt]
   {\tiny Agente D}\\[2pt]
   {\tiny Embeddings duales}\\
   {\tiny Cosine similarity}\\
   {\tiny Verificaci\'on post-ins.}\\
   {\tiny Batch statistics}};

\node[phase, fill=agentE!10, draw=agentE] (p5) at (14,0)
  {\textbf{FASE 5}\\[2pt] Exportaci\'on\\PDF\\[4pt]
   {\tiny Agente E}\\[2pt]
   {\tiny ReportLab}\\
   {\tiny 3 formatos}\\
   {\tiny Histogramas dific.}\\
   {\tiny Agrupado x tema}};

% I/O
\node[io] (in) at (-2.5,0) {Documentos\\legislativos};
\node[io] (out) at (16.5,0) {Ex\'amenes\\PDF};

% Stores
\node[store, fill=agentD!8] (s4) at (10.5,-2.2) {SQLite DB V2};

% Caches
\node[cachenode] (c1) at (0,-2.2) {JSON cache};
\node[cachenode] (c2) at (0,-2.9) {Patterns};
\node[cachenode] (c3) at (3.5,-2.2) {Chunk usage};
\node[cachenode] (c4) at (7,-2.2) {Coordination};

% Flechas principales
\draw[arr, color=pipecolor] (in) -- (p1);
\draw[arr, color=agentZ] (p1) -- node[above, font=\tiny] {chunks} (p2);
\draw[arr, color=agentB] (p2) -- node[above, font=\tiny] {preguntas} (p3);
\draw[arr, color=agentC] (p3) -- node[above, font=\tiny] {validadas} (p4);
\draw[arr, color=agentD] (p4) -- node[above, font=\tiny] {IDs} (p5);
\draw[arr, color=pipecolor] (p5) -- (out);

% Stores arrows
\draw[darr, color=agentZ] (p1) -- (c1);
\draw[darr, color=agentZ] (p1.south) -- ++(0,-0.6) -| (c2);
\draw[darr, color=agentB] (p2) -- (c3);
\draw[darr, color=agentC] (p3) -- (c4);
\draw[darr, color=agentD] (p4) -- (s4);
\draw[darr, color=agentE] (s4.east) -| (p5.south);

% Retry
\draw[arr, color=red!60, line width=0.7pt]
  (p3.south) -- ++(0,-0.7) -| node[below, font=\tiny\itshape, pos=0.25, color=red!60]
  {auto\_fail: retry con feedback (m\'ax 3)} (p2.south);

\end{tikzpicture}
\caption{Flujo de datos detallado del pipeline con capas de cach\'e (dorado).
Cada fase corresponde a un agente especializado. El bucle de retroalimentaci\'on
entre las Fases 2 y 3 garantiza calidad m\'inima antes de la persistencia.}
\label{fig:pipeline}
\end{figure*}

% ============================================================================
\section{Capas de Cach\'e}
\label{sec:cache}
% ============================================================================

El sistema implementa \textbf{4 capas de cach\'e} independientes para optimizar
rendimiento y prevenir p\'erdida de trabajo (Figura~\ref{fig:cache}).

\begin{figure}[H]
\centering
\begin{tikzpicture}[
  cache/.style={
    rectangle, rounded corners=4pt, draw=cachegold, line width=0.8pt,
    minimum width=3.6cm, minimum height=0.7cm,
    text width=3.4cm, align=center, font=\scriptsize, fill=cachegold!8
  },
  desc/.style={font=\tiny\itshape, text=codegray, text width=3.4cm, align=left}
]

% Caches
\node[cache] (c1) at (0,0) {\textbf{Chunk Cache}\\{\tiny input\_docs/rewritten/*.json}};
\node[cache] (c2) at (0,-1.6) {\textbf{Chunk Usage Cache}\\{\tiny tracks processed chunks}};
\node[cache] (c3) at (0,-3.2) {\textbf{Coordination Cache}\\{\tiny Agent Z decisions}};
\node[cache] (c4) at (0,-4.8) {\textbf{Pattern Registry}\\{\tiny patterns.json}};

% Descriptions
\node[desc, right=0.1cm of c1] {Chunks fragmentados como JSON. Evita re-ejecutar Agente Z.
Guardado incremental cada N chunks.};
\node[desc, right=0.1cm of c2] {Registro de qu\'e chunks han sido procesados (count,
last\_used, doc\_id, chunk\_index).};
\node[desc, right=0.1cm of c3] {Decisiones de coordinaci\'on LLM/heur\'istica sobre
l\'imites de fragmentaci\'on.};
\node[desc, right=0.1cm of c4] {Patrones regex aprendidos de limpieza. Persisten entre
documentos y ejecuciones.};

% Agent labels
\node[font=\tiny\bfseries, color=agentZ] at (-2.5,0) {Agente Z};
\node[font=\tiny\bfseries, color=agentB] at (-2.5,-1.6) {Pipeline};
\node[font=\tiny\bfseries, color=agentZ] at (-2.5,-3.2) {Agente Z};
\node[font=\tiny\bfseries, color=agentZ] at (-2.5,-4.8) {Agente Z};

\end{tikzpicture}
\caption{Cuatro capas de cach\'e del sistema. Cada capa previene trabajo
redundante y permite reanudar ejecuciones interrumpidas.}
\label{fig:cache}
\end{figure}

% ============================================================================
\section{Procesamiento Paralelo}
\label{sec:parallel}
% ============================================================================

El sistema implementa procesamiento paralelo a m\'ultiples niveles
(Figura~\ref{fig:parallel}).

\begin{figure}[H]
\centering
\begin{tikzpicture}[
  worker/.style={
    rectangle, rounded corners=3pt, draw=agentB, line width=0.6pt,
    minimum width=1.6cm, minimum height=0.7cm,
    font=\tiny\bfseries, align=center, fill=agentB!10
  },
  queue/.style={
    rectangle, rounded corners=3pt, draw=agentD, line width=0.6pt,
    minimum width=2.2cm, minimum height=0.5cm,
    font=\tiny\bfseries, align=center, fill=agentD!10
  },
  arr/.style={-{Stealth[length=4pt]}, line width=0.6pt}
]

% Thread Pool
\node[rectangle, rounded corners=5pt, draw=agentB, fill=agentB!5,
      line width=0.8pt, font=\scriptsize\bfseries, minimum width=5.5cm,
      minimum height=0.6cm] (pool) at (0,0) {ThreadPoolExecutor};

% Workers
\node[worker] (w1) at (-2,-1) {Worker 1\\B+C};
\node[worker] (w2) at (0,-1) {Worker 2\\B+C};
\node[worker] (w3) at (2,-1) {Worker N\\B+C};
\node[font=\tiny, color=codegray] at (1,-1) {\ldots};

% Queue
\node[queue] (q) at (0,-2.2) {Cola de Escritura};

% DB
\node[rectangle, rounded corners=3pt, draw=dbcolor, fill=dbcolor!10,
      font=\tiny\bfseries, minimum width=2cm, minimum height=0.5cm]
  (db) at (0,-3.2) {ThreadWriter $\to$ SQLite};

% Arrows
\draw[arr, color=agentB] (pool) -- (w1);
\draw[arr, color=agentB] (pool) -- (w2);
\draw[arr, color=agentB] (pool) -- (w3);
\foreach \w in {w1,w2,w3} {
  \draw[arr, color=agentD] (\w) -- (q);
}
\draw[arr, color=dbcolor] (q) -- (db);

% Config
\node[font=\tiny\ttfamily, color=codegray, right=0.2cm of pool]
  {--workers 4};

\end{tikzpicture}
\caption{Arquitectura de procesamiento paralelo. M\'ultiples workers ejecutan
Agentes B+C en paralelo. Un hilo escritor (\texttt{ThreadWriter}) basado en
cola garantiza escrituras at\'omicas a SQLite.}
\label{fig:parallel}
\end{figure}

Caracter\'isticas del sistema paralelo:
\begin{itemize}[nosep,leftmargin=*]
  \item \texttt{ThreadPoolExecutor} con workers configurables (\texttt{--workers N})
  \item Cada worker ejecuta Agentes B y C independientemente por chunk
  \item Cola (\texttt{Queue}) para encolar resultados de escritura
  \item \texttt{ThreadWriter} consume la cola y escribe at\'omicamente a SQLite
  \item Particionado din\'amico de chunks entre workers
  \item Chunk usage cache compartido para evitar reprocesamiento
\end{itemize}

% ============================================================================
\section{Puntos de Entrada}
% ============================================================================

El sistema ofrece cuatro puntos de entrada seg\'un el caso de uso (Figura~\ref{fig:entrypoints}).

\begin{figure}[H]
\centering
\begin{tikzpicture}[
  entry/.style={
    rectangle, rounded corners=4pt, minimum width=3.4cm, minimum height=0.8cm,
    text width=3.2cm, align=center, font=\scriptsize\ttfamily,
    draw, line width=0.8pt, fill=codeblue!10, draw=codeblue
  },
  desc/.style={font=\tiny, text width=3.2cm, align=left},
  agents/.style={
    rectangle, rounded corners=2pt, minimum height=0.5cm,
    font=\tiny\bfseries, draw, line width=0.5pt, inner sep=2pt
  }
]

% Entries
\node[entry] (e1) at (0,0) {main.py};
\node[entry] (e2) at (0,-2.2) {run\_rewriter.py};
\node[entry] (e3) at (0,-4.4) {run\_bcde\_pipeline.py};
\node[entry] (e4) at (0,-6.6) {extract\_and\_generate\_pdf.py};

% Descriptions
\node[desc, right=0.3cm of e1] (d1) {Pipeline completo. Ejecuta Z, B, C, D, E en secuencia.
Tracking con questions\_done.csv.};
\node[desc, right=0.3cm of e2] (d2) {Solo fragmentaci\'on. Ejecuta Agente Z y cachea en JSON.
Auto-topic detection.};
\node[desc, right=0.3cm of e3] (d3) {Pipeline sin chunking. Usa chunks cacheados. Workers
paralelos con ThreadPoolExecutor.};
\node[desc, right=0.3cm of e4] (d4) {Solo exportaci\'on. Genera PDFs desde la base de datos.
3 formatos disponibles.};

% Agent indicators
\node[agents, fill=agentZ!15, right=0.1cm of d1, yshift=0.15cm] {\textcolor{agentZ}{Z}};
\node[agents, fill=agentB!15, right=0.55cm of d1, yshift=0.15cm] {\textcolor{agentB}{B}};
\node[agents, fill=agentC!15, right=1cm of d1, yshift=0.15cm] {\textcolor{agentC}{C}};
\node[agents, fill=agentD!15, right=1.45cm of d1, yshift=0.15cm] {\textcolor{agentD}{D}};
\node[agents, fill=agentE!15, right=1.9cm of d1, yshift=0.15cm] {\textcolor{agentE}{E}};

\node[agents, fill=agentZ!15, right=0.1cm of d2, yshift=0.15cm] {\textcolor{agentZ}{Z}};

\node[agents, fill=agentB!15, right=0.1cm of d3, yshift=0.15cm] {\textcolor{agentB}{B}};
\node[agents, fill=agentC!15, right=0.55cm of d3, yshift=0.15cm] {\textcolor{agentC}{C}};
\node[agents, fill=agentD!15, right=1cm of d3, yshift=0.15cm] {\textcolor{agentD}{D}};
\node[agents, fill=agentE!15, right=1.45cm of d3, yshift=0.15cm] {\textcolor{agentE}{E}};

\node[agents, fill=agentE!15, right=0.1cm of d4, yshift=0.15cm] {\textcolor{agentE}{E}};

\end{tikzpicture}
\caption{Puntos de entrada del sistema y agentes que activan.}
\label{fig:entrypoints}
\end{figure}

\subsection{main.py --- Pipeline Completo}

Orquesta la ejecuci\'on secuencial de los cinco agentes. Es el punto de entrada
principal para producci\'on. Utiliza \texttt{questions\_done.csv} para tracking
de documentos ya procesados.

\begin{tcolorbox}[colback=bglight, colframe=bordercolor, boxrule=0.5pt, arc=2pt]
{\small\ttfamily
python main.py \textbackslash\\
\hspace*{1em}--doc "Tema-18" \textbackslash\\
\hspace*{1em}--questions 5 \textbackslash\\
\hspace*{1em}--topic 101 \textbackslash\\
\hspace*{1em}--academy 1 \textbackslash\\
\hspace*{1em}--skip-agent-c ~~\# opcional
}
\end{tcolorbox}

\subsection{run\_rewriter.py --- Solo Agente Z}

Ejecuta \'unicamente la fase de fragmentaci\'on con \textbf{auto-topic detection}
y l\'imites configurables.

\begin{tcolorbox}[colback=bglight, colframe=bordercolor, boxrule=0.5pt, arc=2pt]
{\small\ttfamily
python run\_rewriter.py \textbackslash\\
\hspace*{1em}--auto-topic \textbackslash\\
\hspace*{1em}--limit 5 \textbackslash\\
\hspace*{1em}--limit-chunks 10 \textbackslash\\
\hspace*{1em}--force ~~\# re-procesar existentes
}
\end{tcolorbox}

\subsection{run\_bcde\_pipeline.py --- Pipeline Paralelo}

El punto de entrada m\'as flexible. Consume chunks cacheados y ejecuta B, C, D, E
con \textbf{procesamiento paralelo}.

\begin{tcolorbox}[colback=bglight, colframe=bordercolor, boxrule=0.5pt, arc=2pt]
{\small\ttfamily
python run\_bcde\_pipeline.py \textbackslash\\
\hspace*{1em}--doc "Tema-18" \textbackslash\\
\hspace*{1em}--questions-per-chunk 5 \textbackslash\\
\hspace*{1em}--workers 4 \textbackslash\\
\hspace*{1em}--skip-agent-c
}
\end{tcolorbox}

\subsection{extract\_and\_generate\_pdf.py --- Solo PDFs}

Utilidad para generar PDFs a partir de preguntas ya existentes en la base
de datos. No ejecuta ning\'un agente generativo.

% ============================================================================
\section{Base de Datos Local (Esquema V2)}
% ============================================================================

El sistema utiliza \textbf{SQLite} (motor de base de datos relacional embebido,
sin necesidad de servidor) como base de datos local, con un esquema V2
dise\~nado para ser compatible con Supabase para futuras migraciones a la nube.

\subsection{Esquema Relacional}

La Figura~\ref{fig:database} muestra el esquema Entity-Relationship de la base
de datos.

\begin{figure}[H]
\centering
\begin{tikzpicture}[
  header/.style={
    rectangle, minimum width=3.6cm, minimum height=0.45cm,
    font=\scriptsize\bfseries, align=center, text=white
  },
  field/.style={font=\tiny\ttfamily, anchor=west},
  pk/.style={font=\tiny\ttfamily\bfseries, anchor=west, text=codeblue},
  fk/.style={font=\tiny\ttfamily\itshape, anchor=west, text=agentD},
  groupl/.style={font=\tiny\bfseries\itshape, anchor=west, text=codegray}
]

% --- Table: questions ---
\node[header, fill=dbcolor] (qh) at (0,0) {questions};
\draw[line width=0.8pt] (-1.8,-0.23) rectangle (1.8,-5.83);
\node[pk] at (-1.7,-0.5) {PK id INTEGER};
\node[groupl] at (-1.7,-0.8) {-- Contenido --};
\node[field] at (-1.7,-1.1) {question TEXT};
\node[field] at (-1.7,-1.4) {tip TEXT};
\node[field] at (-1.7,-1.7) {article TEXT};
\node[field] at (-1.7,-2.0) {topic INTEGER};
\node[field] at (-1.7,-2.3) {academy\_id INTEGER};
\node[groupl] at (-1.7,-2.6) {-- Calidad --};
\node[field] at (-1.7,-2.9) {faithfulness\_score REAL};
\node[field] at (-1.7,-3.2) {relevancy\_score REAL};
\node[field] at (-1.7,-3.5) {needs\_manual\_review BOOL};
\node[groupl] at (-1.7,-3.8) {-- Deduplicaci\'on --};
\node[field] at (-1.7,-4.1) {is\_duplicate BOOLEAN};
\node[field] at (-1.7,-4.4) {duplicate\_of INTEGER};
\node[groupl] at (-1.7,-4.7) {-- Trazabilidad --};
\node[field] at (-1.7,-5.0) {source\_document, source\_page};
\node[field] at (-1.7,-5.3) {llm\_model, retry\_count};
\node[field] at (-1.7,-5.6) {distractor\_techniques JSON};

% --- Table: question_options ---
\node[header, fill=agentB!80!black] (oh) at (6,0) {question\_options};
\draw[line width=0.8pt] (4.2,-0.23) rectangle (7.8,-2.33);
\node[pk] at (4.3,-0.5) {PK id INTEGER};
\node[fk] at (4.3,-0.8) {FK question\_id INTEGER};
\node[field] at (4.3,-1.1) {answer TEXT};
\node[field] at (4.3,-1.4) {is\_correct BOOLEAN};
\node[field] at (4.3,-1.7) {option\_order INTEGER};
\node[field] at (4.3,-2.0) {created\_at TIMESTAMP};

% --- Table: batches ---
\node[header, fill=agentD!80!black] (bh) at (6,-3.2) {batches};
\draw[line width=0.8pt] (4.2,-3.43) rectangle (7.8,-6.13);
\node[pk] at (4.3,-3.7) {PK id INTEGER};
\node[field] at (4.3,-4.0) {batch\_name TEXT};
\node[field] at (4.3,-4.3) {topic INTEGER};
\node[field] at (4.3,-4.6) {total\_questions INTEGER};
\node[field] at (4.3,-4.9) {avg\_faithfulness REAL};
\node[field] at (4.3,-5.2) {avg\_relevancy REAL};
\node[field] at (4.3,-5.5) {source\_documents JSON};
\node[field] at (4.3,-5.8) {status TEXT};

% Relationships
\draw[-{Stealth[length=5pt]}, line width=0.8pt, color=agentD]
  (1.8,-0.8) -- node[above, font=\tiny] {1:N} (4.2,-0.8);

% Views label
\node[rectangle, rounded corners=3pt, draw=codeblue, fill=codeblue!8,
      font=\tiny\bfseries, text width=3cm, align=center] at (6,-7)
  {Vistas SQL:\\questions\_with\_options\\questions\_full\\batch\_quality\_stats};

\end{tikzpicture}
\caption{Esquema Entity-Relationship de la base de datos SQLite V2.
Incluye vistas SQL precalculadas para consultas frecuentes.}
\label{fig:database}
\end{figure}

\subsection{Modelo Question V2}

El modelo \texttt{QuestionV2} (\texttt{models/question\_v2.py}) es un modelo Pydantic
con m\'as de 30 campos organizados en grupos:

\begin{table}[H]
\centering
\caption{Grupos de campos del modelo Question V2}
\label{tab:questionv2}
\scriptsize
\begin{tabularx}{\columnwidth}{lX}
\toprule
\textbf{Grupo} & \textbf{Campos principales} \\
\midrule
Contenido & question, tip, article, topic \\
Calidad & faithfulness\_score, relevancy\_score \\
Dedup & is\_duplicate, duplicate\_of \\
Trazabilidad & source\_chunk\_id, source\_document, source\_page \\
LLM & llm\_model, by\_llm, generation\_method \\
Retry & retry\_count, needs\_manual\_review, manual\_review\_reason \\
Multimedia & question\_image\_url, retro\_image\_url, retro\_audio\_* \\
Estad\'isticas & num\_answered, num\_fails, num\_empty \\
Desaf\'io & challenge\_by\_tutor, challenge\_reason \\
\bottomrule
\end{tabularx}
\end{table}

M\'etodos computados del modelo:
\begin{itemize}[nosep,leftmargin=*]
  \item \texttt{calculate\_difficult\_rate()}: Tasa de dificultad basada en estad\'isticas
  \item \texttt{is\_high\_quality()}: Evaluaci\'on r\'apida de calidad
  \item \texttt{should\_retry()}: Determina si necesita regeneraci\'on
  \item \texttt{get\_text\_for\_embedding()}: Texto optimizado para embeddings
  \item \texttt{from\_legacy\_dict()} / \texttt{to\_legacy\_dict()}: Compatibilidad V1
\end{itemize}

\subsection{Vistas SQL}

El esquema incluye 3 vistas precalculadas:
\begin{itemize}[nosep,leftmargin=*]
  \item \texttt{questions\_with\_options}: JOIN pregunta + opciones
  \item \texttt{questions\_full}: Vista completa con todos los metadatos
  \item \texttt{batch\_quality\_stats}: Estad\'isticas de calidad por batch
\end{itemize}

\subsection{Repositorio de Acceso}

La clase \texttt{QuestionRepository} implementa el patr\'on Repository (patr\'on de
dise\~no que abstrae el acceso a datos) con context manager (gestor de contexto
que asegura apertura y cierre autom\'atico de conexiones) para conexiones seguras:

\begin{tcolorbox}[colback=bglight, colframe=bordercolor, boxrule=0.5pt, arc=2pt]
{\small\ttfamily
repo = QuestionRepository(db\_path)\\
\# Insercion individual\\
qid = repo.insert(question\_data)\\
\# Insercion por lotes\\
ids = repo.insert\_batch(questions)\\
\# Consultas\\
q = repo.get\_by\_id(qid)\\
all\_q = repo.get\_all(topic=101)\\
exists = repo.exists(question\_id)\\
stats = repo.get\_stats()
}
\end{tcolorbox}

% ============================================================================
\section{Estado del Grafo (GraphState)}
% ============================================================================

El estado compartido entre agentes se define como un \texttt{TypedDict}
(diccionario tipado de Python que define la estructura de datos esperada) en
\texttt{graph/state.py} con aproximadamente 120 campos organizados en 7 grupos
(Figura~\ref{fig:state}).

\begin{figure}[H]
\centering
\begin{tikzpicture}[
  group/.style={
    rectangle, rounded corners=3pt, draw, line width=0.6pt,
    text width=3.4cm, minimum height=0.7cm, align=center,
    font=\scriptsize, inner sep=4pt
  },
  detail/.style={font=\tiny\ttfamily, color=codegray, text width=3cm, align=left},
  arr/.style={-{Stealth[length=4pt]}, line width=0.6pt, color=pipecolor}
]

\node[group, fill=agentZ!10, draw=agentZ] (g1) at (0,0) {\textbf{Input}\\{\tiny docs, metadata, file\_type}};
\node[group, fill=agentZ!10, draw=agentZ] (g2) at (0,-1.2) {\textbf{Chunks}\\{\tiny List[Chunk], substantive\_count}};
\node[group, fill=agentB!10, draw=agentB] (g3) at (0,-2.4) {\textbf{Generaci\'on}\\{\tiny Dict[chunk\_id, List[Q]]}};
\node[group, fill=agentC!10, draw=agentC] (g4) at (0,-3.6) {\textbf{Validaci\'on}\\{\tiny validated, failed, review}};
\node[group, fill=agentC!10, draw=agentC] (g5) at (0,-4.8) {\textbf{Retry}\\{\tiny counts, max\_retries, feedback}};
\node[group, fill=agentD!10, draw=agentD] (g6) at (0,-6.0) {\textbf{Persistencia}\\{\tiny IDs, dedup\_stats, verify\_errors}};
\node[group, fill=pipecolor!10, draw=pipecolor] (g7) at (0,-7.2) {\textbf{Progreso}\\{\tiny step, completed, totals, errors}};

\foreach \i [evaluate=\i as \j using int(\i+1)] in {1,...,6} {
  \draw[arr] (g\i) -- (g\j);
}

% Multi-chunk stats
\node[rectangle, rounded corners=3pt, draw=reactblue, fill=reactblue!8,
      font=\tiny, text width=2.8cm, align=center, line width=0.5pt]
  at (4.5,-2.4) {\textbf{Multi-Chunk Stats}\\enriched\_contexts\\multi\_context\_questions\\chunk\_retriever\_initialized};

\end{tikzpicture}
\caption{Grupos de campos en el \texttt{GraphState}. El estado fluye
secuencialmente a trav\'es de los agentes. Se incluyen estad\'isticas
de multi-chunk y tracking de errores.}
\label{fig:state}
\end{figure}

% ============================================================================
\section{Sistema de Configuraci\'on}
% ============================================================================

La configuraci\'on utiliza \textbf{Pydantic BaseSettings} con un sistema jer\'arquico
de \texttt{ConfigBundle} que agrega modelos por agente (Figura~\ref{fig:config}).

\begin{figure}[H]
\centering
\begin{tikzpicture}[
  cfg/.style={
    rectangle, rounded corners=3pt, draw, line width=0.6pt,
    minimum width=3.4cm, minimum height=0.6cm,
    font=\scriptsize, align=center, fill=bglight
  },
  main/.style={
    rectangle, rounded corners=5pt, draw=codeblue, line width=1pt,
    minimum width=3.8cm, minimum height=0.8cm,
    font=\small\bfseries, align=center, fill=codeblue!10
  },
  param/.style={font=\tiny\ttfamily, color=codegray}
]

\node[main] (settings) at (0,0) {ConfigBundle};

\node[cfg] (openai) at (-2.2,-1.2) {OpenAI Config};
\node[cfg] (agentz) at (2.2,-1.2) {Agent Z Config};
\node[cfg] (agentb) at (-2.2,-2.4) {Agent B Config};
\node[cfg] (agentc) at (2.2,-2.4) {Agent C Config};
\node[cfg] (embed) at (-2.2,-3.6) {Embeddings Config};
\node[cfg] (dedup) at (2.2,-3.6) {Dedup Config};
\node[cfg] (paths) at (0,-4.6) {Paths Config};

\foreach \n in {openai,agentz,agentb,agentc,embed,dedup,paths} {
  \draw[-{Stealth[length=4pt]}, line width=0.5pt] (settings) -- (\n);
}

\node[param, below=0.05cm of openai] {model, api\_key, reasoning\_effort};
\node[param, below=0.05cm of agentz] {chunk\_size, overlap, coord\_mode};
\node[param, below=0.05cm of agentb] {questions\_per\_chunk, max\_context\_tokens};
\node[param, below=0.05cm of agentc] {thresholds, max\_retries, fast\_mode};
\node[param, below=0.05cm of embed] {provider, model, gpu\_support};
\node[param, below=0.05cm of dedup] {similarity: 0.85, batch\_size};
\node[param, below=0.05cm of paths] {input\_dir, output\_dir, db\_path};

% .env
\node[rectangle, rounded corners=3pt, draw=cachegold, fill=cachegold!10,
      font=\scriptsize\bfseries, minimum width=1.6cm] at (4.5,0) {.env};
\draw[-{Stealth[length=4pt]}, line width=0.5pt, color=cachegold] (4.5,0) -- (settings);

\end{tikzpicture}
\caption{Sistema de configuraci\'on jer\'arquico basado en \texttt{ConfigBundle}.
Las variables de entorno se cargan desde \texttt{.env} v\'ia Pydantic BaseSettings.}
\label{fig:config}
\end{figure}

\subsection{Modelos LLM Configurables}

El sistema soporta \textbf{m\'ultiples modelos LLM} para diferentes tareas:

\begin{table}[H]
\centering
\caption{Modelos LLM por tarea}
\label{tab:models}
\scriptsize
\begin{tabularx}{\columnwidth}{lX}
\toprule
\textbf{Tarea} & \textbf{Configuraci\'on} \\
\midrule
Generaci\'on (B) & \texttt{OPENAI\_MODEL} (GPT-5-mini) \\
Evaluaci\'on (C) & \texttt{OPENAI\_MODEL} (mismo u otro) \\
Reescritura (Z) & Modelo dedicado opcional \\
An\'alisis dificultad & \texttt{difficulty\_model} (opcional) \\
Reasoning effort & \texttt{low|medium|high} (GPT-5) \\
\bottomrule
\end{tabularx}
\end{table}

\subsection{Timeouts Granulares}

\begin{table}[H]
\centering
\caption{Configuraci\'on de timeouts por operaci\'on}
\label{tab:timeouts}
\scriptsize
\begin{tabularx}{\columnwidth}{lX}
\toprule
\textbf{Operaci\'on} & \textbf{Timeout (seg)} \\
\midrule
Agente Z (por documento) & 300 \\
Llamadas LLM generales & 120 \\
Agente C (evaluaci\'on) & 120 \\
Operaciones de embedding & 60 \\
\bottomrule
\end{tabularx}
\end{table}

\subsection{Perfiles de Configuraci\'on}

\begin{table}[H]
\centering
\caption{Perfiles de configuraci\'on recomendados}
\label{tab:profiles}
\scriptsize
\begin{tabularx}{\columnwidth}{lXXX}
\toprule
\textbf{Par\'ametro} & \textbf{R\'apido} & \textbf{Balanceado} & \textbf{Estricto} \\
\midrule
faith. threshold & 0.80 & 0.85 & 0.90 \\
relev. threshold & 0.80 & 0.85 & 0.90 \\
questions/chunk & 2 & 5 & 3 \\
max retries & 1 & 3 & 3 \\
similarity dedup & 0.80 & 0.85 & 0.90 \\
skip Agent C & S\'i & No & No \\
fast mode C & S\'i & No & No \\
tip\_min\_words & 30 & 40 & 60 \\
tip\_max\_words & 200 & 200 & 180 \\
\bottomrule
\end{tabularx}
\end{table}

% ============================================================================
\section{Ingenier\'ia de Prompts}
\label{sec:prompts}
% ============================================================================

El sistema utiliza m\'as de \textbf{30 archivos de prompts} organizados en
\texttt{config/prompts/}, gestionados por el \texttt{PromptLoader} con soporte
de templates (Figura~\ref{fig:prompts}).

\begin{figure}[H]
\centering
\begin{tikzpicture}[
  file/.style={
    rectangle, rounded corners=2pt, draw, line width=0.5pt,
    minimum width=2.8cm, minimum height=0.4cm,
    font=\tiny\ttfamily, align=center, fill=bglight
  },
  group/.style={
    rectangle, rounded corners=4pt, draw, line width=0.8pt,
    minimum width=3.2cm, inner sep=4pt
  }
]

% Agent Z
\node[group, fill=agentZ!5, draw=agentZ] (gz) at (0,0) {};
\node[font=\scriptsize\bfseries, color=agentZ] at (0,0.6) {Agente Z};
\node[file] at (0,0.15) {agent\_z\_rewriter\_system.txt};
\node[file] at (0,-0.25) {agent\_z\_rewriter\_user.txt};

% Agent B
\node[group, fill=agentB!5, draw=agentB] (gb) at (0,-1.6) {};
\node[font=\scriptsize\bfseries, color=agentB] at (0,-0.9) {Agente B};
\node[file] at (0,-1.25) {agent\_b\_generation.json};
\node[file] at (0,-1.65) {agent\_b\_simple\_system.txt};
\node[file] at (0,-2.05) {agent\_b\_evaluate\_system.txt};

% Agent C
\node[group, fill=agentC!5, draw=agentC] (gc) at (0,-3.8) {};
\node[font=\scriptsize\bfseries, color=agentC] at (0,-2.8) {Agente C (5 prompts)};
\node[file] at (0,-3.15) {agent\_c\_specialized\_*.txt};
\node[file] at (0,-3.55) {agent\_c\_tip\_consistency\_*.txt};
\node[file] at (0,-3.95) {agent\_c\_tip\_supports\_*.txt};
\node[file] at (0,-4.35) {agent\_c\_validate\_distractor\_*.txt};
\node[file] at (0,-4.75) {agent\_c\_quick\_*.txt};

% Loader
\node[rectangle, rounded corners=4pt, draw=codeblue, fill=codeblue!10,
      line width=0.8pt, font=\scriptsize\bfseries, minimum width=2.4cm,
      minimum height=0.5cm] at (4,-2) {PromptLoader};

\draw[-{Stealth[length=4pt]}, line width=0.5pt, color=codeblue]
  (1.6,-0.1) -- (2.8,-1.8);
\draw[-{Stealth[length=4pt]}, line width=0.5pt, color=codeblue]
  (1.6,-1.6) -- (2.8,-2);
\draw[-{Stealth[length=4pt]}, line width=0.5pt, color=codeblue]
  (1.6,-3.5) -- (2.8,-2.2);

\end{tikzpicture}
\caption{Organizaci\'on de los 30+ archivos de prompts. El \texttt{PromptLoader}
gestiona la carga con soporte de templates y ejemplos few-shot en JSON.}
\label{fig:prompts}
\end{figure}

Los prompts del Agente~B incluyen:
\begin{itemize}[nosep,leftmargin=*]
  \item \textbf{Few-shot examples}: 6 preguntas reales en JSON estructurado con
        pregunta, opciones, respuesta correcta, tip y t\'ecnicas de distracci\'on
  \item \textbf{System prompt}: Instrucciones detalladas sobre formato, estilo,
        nivel de dificultad y t\'ecnicas de distracci\'on permitidas
  \item \textbf{Modo simple}: Prompt reducido para iteraci\'on r\'apida
  \item \textbf{Evaluate prompt}: Para el paso de evaluaci\'on guardian pre-generaci\'on
\end{itemize}

% ============================================================================
\section{Stack Tecnol\'ogico}
% ============================================================================

\begin{figure}[H]
\centering
\begin{tikzpicture}[
  layer/.style={
    rectangle, rounded corners=4pt, minimum width=7.2cm, minimum height=0.9cm,
    text width=7cm, align=center, font=\scriptsize, draw, line width=0.6pt
  }
]

\node[layer, fill=agentE!10] (l1) at (0,0)
  {\textbf{Presentaci\'on}: ReportLab (PDF), Matplotlib (histogramas)};
\node[layer, fill=agentB!10] (l2) at (0,-1.1)
  {\textbf{Agentes}: LangGraph (grafos de agentes), StateGraph, ReAct, create\_react\_agent};
\node[layer, fill=agentC!10] (l3) at (0,-2.2)
  {\textbf{Calidad}: Ragas (faithfulness, relevancy), Eval. Especializada};
\node[layer, fill=agentZ!10] (l4) at (0,-3.3)
  {\textbf{LLM} (modelos de lenguaje): OpenAI GPT-5-mini, FAISS, sentence-transformers};
\node[layer, fill=agentD!10] (l5) at (0,-4.4)
  {\textbf{Datos}: SQLite, Pydantic V2, pypdf, pdfplumber, BaseSettings};
\node[layer, fill=pipecolor!10] (l6) at (0,-5.5)
  {\textbf{Infra}: Python 3.11+, ThreadPoolExecutor, tenacity (reintentos), tqdm (barras de progreso), dotenv};

\end{tikzpicture}
\caption{Stack tecnol\'ogico organizado por capas.}
\label{fig:stack}
\end{figure}

% ============================================================================
\section{Gesti\'on de Errores y Robustez}
% ============================================================================

El sistema implementa m\'ultiples mecanismos de resiliencia:

\begin{figure}[H]
\centering
\begin{tikzpicture}[
  mechanism/.style={
    rectangle, rounded corners=3pt, draw, line width=0.6pt,
    minimum width=3.6cm, minimum height=0.55cm,
    text width=3.4cm, align=center, font=\scriptsize
  },
  arr/.style={-{Stealth[length=4pt]}, line width=0.5pt}
]

\node[mechanism, fill=agentE!10, draw=agentE] (r1) at (0,0)
  {\textbf{Retry autom\'atico}\\{\tiny max 3/chunk + 2/pregunta}};
\node[mechanism, fill=agentC!10, draw=agentC] (r2) at (0,-1)
  {\textbf{Retroalimentaci\'on espec\'ifica}\\{\tiny motivo + criterios fallidos}};
\node[mechanism, fill=cachegold!10, draw=cachegold] (r3) at (0,-2)
  {\textbf{Guardado incremental}\\{\tiny cada N chunks a JSON}};
\node[mechanism, fill=agentD!10, draw=agentD] (r4) at (0,-3)
  {\textbf{Verificaci\'on post-inserci\'on}\\{\tiny \_verify\_inserted()}};
\node[mechanism, fill=agentZ!10, draw=agentZ] (r5) at (0,-4)
  {\textbf{Timeout por operaci\'on}\\{\tiny doc: 300s, LLM: 120s}};
\node[mechanism, fill=pipecolor!10, draw=pipecolor] (r6) at (0,-5)
  {\textbf{Tracking de errores}\\{\tiny GraphState.errors list}};
\node[mechanism, fill=codegreen!10, draw=codegreen] (r7) at (0,-6)
  {\textbf{Degradaci\'on elegante}\\{\tiny fast mode, skip Agent C}};

\end{tikzpicture}
\caption{Siete mecanismos de resiliencia y robustez del sistema.}
\label{fig:resilience}
\end{figure}

El tracking de errores en el \texttt{GraphState} almacena cada error como un
diccionario con \texttt{\{node, chunk\_id, error, timestamp\}}, permitiendo
diagn\'ostico post-ejecuci\'on.

% ============================================================================
\section{Flujos de Trabajo Recomendados}
% ============================================================================

\subsection{Desarrollo / Prueba R\'apida}

\begin{tcolorbox}[colback=bglight, colframe=bordercolor, boxrule=0.5pt, arc=2pt]
{\small\ttfamily
\# Paso 1: Fragmentar 1 documento, max 5 chunks\\
python run\_rewriter.py --limit 1 --limit-chunks 5\\[4pt]
\# Paso 2: Generar sin validacion (fast)\\
python run\_bcde\_pipeline.py --questions-per-chunk 2 --skip-agent-c
}
\end{tcolorbox}

\subsection{Producci\'on}

\begin{tcolorbox}[colback=bglight, colframe=bordercolor, boxrule=0.5pt, arc=2pt]
{\small\ttfamily
\# Pipeline completo con validacion dual\\
python main.py --questions 5 --topic 101 --academy 1
}
\end{tcolorbox}

\subsection{Iteraci\'on Paralela sobre Documento}

\begin{tcolorbox}[colback=bglight, colframe=bordercolor, boxrule=0.5pt, arc=2pt]
{\small\ttfamily
\# Usar chunks cacheados, 4 workers paralelos\\
python run\_bcde\_pipeline.py --doc "Tema-18" \textbackslash\\
\hspace*{1em}--questions-per-chunk 5 --workers 4
}
\end{tcolorbox}

% ============================================================================
\section{M\'etricas y Trazabilidad}
% ============================================================================

El sistema mantiene trazabilidad completa en cada pregunta generada:

\begin{itemize}[nosep,leftmargin=*]
  \item \textbf{Origen}: documento fuente, p\'agina, chunk ID, chunks relacionados
  \item \textbf{Calidad}: scores de faithfulness y relevancy, evaluaci\'on especializada
  \item \textbf{Generaci\'on}: modelo LLM, prompt usado, m\'etodo, context\_strategy,
        generation\_time, reasoning\_summary
  \item \textbf{Deduplicaci\'on}: flag de duplicado, ID del original, similarity score
  \item \textbf{T\'ecnicas}: array JSON de t\'ecnicas de distracci\'on usadas
  \item \textbf{Reintentos}: contador de intentos, flag de revisi\'on manual, raz\'on
  \item \textbf{Multimedia}: URLs de im\'agenes y audio para retroalimentaci\'on
  \item \textbf{Estad\'isticas de uso}: num\_answered, num\_fails, num\_empty
\end{itemize}

Esta trazabilidad permite:
\begin{enumerate}[nosep,leftmargin=*]
  \item Auditar la calidad de las preguntas generadas
  \item Identificar chunks que producen preguntas de baja calidad
  \item Optimizar prompts bas\'andose en m\'etricas hist\'oricas
  \item Detectar t\'ecnicas de distracci\'on m\'as efectivas
  \item Migrar a Supabase manteniendo toda la metadata
\end{enumerate}

% ============================================================================
\section{Conclusiones}
% ============================================================================

El sistema multi-agente presentado ofrece una soluci\'on completa y extensible
para la generaci\'on autom\'atica de preguntas tipo test. La separaci\'on en cinco
agentes especializados, cada uno implementado como StateGraph con patr\'on ReAct,
permite:

\begin{itemize}[nosep,leftmargin=*]
  \item \textbf{Modularidad}: Cada agente puede evolucionar independientemente,
        con sus propios prompts, configuraci\'on y herramientas
  \item \textbf{Calidad}: El sistema multi-estrategia de evaluaci\'on (5 evaluaciones
        en Agente~C) con retry autom\'atico garantiza est\'andares m\'inimos
  \item \textbf{Escalabilidad}: Procesamiento paralelo con \texttt{ThreadPoolExecutor},
        4 capas de cach\'e, y cola de escritura at\'omica
  \item \textbf{Inteligencia}: Recuperaci\'on multi-chunk con FAISS, aprendizaje
        de patrones, y embeddings duales
  \item \textbf{Trazabilidad}: Cada pregunta conserva 30+ campos de metadatos
        incluyendo origen, calidad, modelo, t\'ecnicas y estad\'isticas
  \item \textbf{Robustez}: 7 mecanismos de resiliencia incluyendo timeouts,
        retries, verificaci\'on y degradaci\'on elegante
  \item \textbf{Flexibilidad}: 4 puntos de entrada, 3 perfiles de configuraci\'on,
        y modo fast para iteraci\'on r\'apida
\end{itemize}

El dise\~no del esquema de base de datos V2 compatible con Supabase, con vistas SQL
precalculadas y modelo Question V2 con compatibilidad legacy, facilita la
futura migraci\'on a infraestructura cloud sin p\'erdida de datos ni metadata.

% ============================================================================
\appendix
\section{Estructura de Archivos}
% ============================================================================

\begin{tcolorbox}[colback=bglight, colframe=bordercolor, boxrule=0.5pt, arc=2pt]
{\scriptsize\ttfamily
multiagent\_question\_generation/\\
\hspace*{0.5em}agents/\\
\hspace*{1.5em}agent\_z\_rewriter.py ~~~~~~(2584 LOC)\\
\hspace*{1.5em}agent\_b\_generator.py ~~~~~(2143 LOC)\\
\hspace*{1.5em}agent\_c\_evaluator.py ~~~~~(1751 LOC)\\
\hspace*{1.5em}agent\_d\_persistence.py ~~~(344 LOC)\\
\hspace*{1.5em}agent\_e\_pdf\_generator.py ~(1159 LOC)\\
\hspace*{0.5em}config/\\
\hspace*{1.5em}settings.py\\
\hspace*{1.5em}config\_models.py ~~(ConfigBundle)\\
\hspace*{1.5em}thresholds.py\\
\hspace*{1.5em}prompts/ ~~~~(30+ archivos .txt/.json)\\
\hspace*{0.5em}database/\\
\hspace*{1.5em}questions.db\\
\hspace*{1.5em}init\_schema\_v2.sql ~(3 vistas SQL)\\
\hspace*{1.5em}repository.py\\
\hspace*{0.5em}graph/\\
\hspace*{1.5em}state.py ~~~~(\~{}120 campos GraphState)\\
\hspace*{0.5em}models/\\
\hspace*{1.5em}question.py, question\_v2.py ~(30+ campos)\\
\hspace*{1.5em}question\_option.py\\
\hspace*{1.5em}chunk.py ~~~~(coherence analysis)\\
\hspace*{1.5em}quality\_metrics.py\\
\hspace*{0.5em}services/\\
\hspace*{1.5em}chunk\_retriever.py ~(FAISS multi-chunk)\\
\hspace*{0.5em}utils/\\
\hspace*{1.5em}llm\_factory.py\\
\hspace*{1.5em}loaders.py ~~(PDF/TXT/MD)\\
\hspace*{1.5em}embeddings.py ~~(dual provider)\\
\hspace*{1.5em}metadata\_filter.py\\
\hspace*{1.5em}prompt\_loader.py ~~(templates)\\
\hspace*{1.5em}pdf\_visualizer.py\\
\hspace*{1.5em}rewrite\_pattern\_registry.py\\
\hspace*{0.5em}main.py\\
\hspace*{0.5em}run\_bcde\_pipeline.py ~(ThreadPoolExecutor)\\
\hspace*{0.5em}run\_rewriter.py\\
\hspace*{0.5em}extract\_and\_generate\_pdf.py
}
\end{tcolorbox}

% ============================================================================
\section{Ejemplo Real: Chunk Fragmentado}
\label{app:chunk}
% ============================================================================

A continuaci\'on se muestra un chunk real generado por el Agente~Z a partir del
Tema~21 (Ley Org\'anica 1/2004 de Violencia de G\'enero):

\begin{tcolorbox}[colback=agentZ!3, colframe=agentZ, title={\small Chunk real --- Art\'iculo 1 (L.O. 1/2004)},
  fonttitle=\bfseries\small, boxrule=0.5pt, arc=3pt, breakable]
{\scriptsize
\textbf{chunk\_id:} \texttt{Tema-21\_coherent\_13740b8e}\\
\textbf{source:} Tema-21-Ingreso-a-Guardia-Civil\\
\textbf{token\_count:} 301 ~|~ \textbf{chunk\_type:} partial\_content ~|~ \textbf{law\_reference:} LEY ORG\'ANICA 1/2004\\[4pt]
\rule{\linewidth}{0.3pt}\\[2pt]
Ley referida: LEY ORG\'ANICA 1/2004

\textbf{Art\'iculo 1. Objeto de la Ley.}
1. La presente Ley tiene por objeto actuar contra la violencia que, como manifestaci\'on
de la discriminaci\'on, la situaci\'on de desigualdad y las relaciones de poder de los
hombres sobre las mujeres, se ejerce sobre \'estas por parte de quienes sean o hayan
sido sus c\'onyuges o de quienes est\'en o hayan estado ligados a ellas por relaciones
similares de afectividad, aun sin convivencia.

2. Por esta ley se establecen medidas de protecci\'on integral cuya finalidad es prevenir,
sancionar y erradicar esta violencia y prestar asistencia a las mujeres, a sus hijos
menores y a los menores sujetos a su tutela, o guarda y custodia, v\'ictimas de esta violencia.

3. La violencia de g\'enero [\ldots] comprende todo acto de violencia f\'isica y psicol\'ogica,
incluidas las agresiones a la libertad sexual, las amenazas, las coacciones o la privaci\'on
arbitraria de libertad.
}
\end{tcolorbox}

Metadatos asociados al chunk:

\begin{tcolorbox}[colback=bglight, colframe=bordercolor, boxrule=0.5pt, arc=2pt]
{\scriptsize\ttfamily
\{\\
\hspace*{1em}"topic": 21,\\
\hspace*{1em}"chunk\_index": 1,\\
\hspace*{1em}"total\_chunks": 79,\\
\hspace*{1em}"coherent": true,\\
\hspace*{1em}"chunk\_type": "partial\_content",\\
\hspace*{1em}"rewritten": false,\\
\hspace*{1em}"law\_reference": "LEY ORG\'ANICA 1/2004"\\
\}
}
\end{tcolorbox}

% ============================================================================
\section{Ejemplo Real: Pregunta Generada}
\label{app:question}
% ============================================================================

El siguiente ejemplo muestra una pregunta real generada por el Agente~B,
evaluada por el Agente~C, y almacenada por el Agente~D:

\begin{tcolorbox}[colback=agentB!3, colframe=agentB,
  title={\small Pregunta generada --- C\'odigo Penal, Art. 138},
  fonttitle=\bfseries\small, boxrule=0.5pt, arc=3pt]
{\small
\textbf{Pregunta:} \textquestiondown Cu\'al es la pena por homicidio seg\'un el art\'iculo 138
del C\'odigo Penal?\\[4pt]
\begin{tabular}{cl}
\textbf{A)} & Prisi\'on de 10 a 15 a\~nos \quad \checkmark \\
\textbf{B)} & Prisi\'on de 5 a 10 a\~nos \\
\textbf{C)} & Prisi\'on de 15 a 20 a\~nos \\
\textbf{D)} & Prisi\'on de 20 a 25 a\~nos \\
\end{tabular}
}
\end{tcolorbox}

\begin{tcolorbox}[colback=agentC!3, colframe=agentC,
  title={\small Metadatos de calidad y trazabilidad},
  fonttitle=\bfseries\small, boxrule=0.5pt, arc=3pt]
{\scriptsize\ttfamily
\textbf{Tip (explicaci\'on):} El art\'iculo 138 CP establece\\
que el que matare a otro ser\'a castigado, como reo\\
de homicidio, con pena de prisi\'on de 10 a 15 a\~nos.\\[4pt]
\textbf{Art\'iculo fuente:} Art. 138.1 del C\'odigo Penal\\[4pt]
\textbf{Faithfulness:} 0.95 \quad \textbf{Relevancy:} 0.90\\
\textbf{Clasificaci\'on:} AUTO\_PASS\\
\textbf{T\'ecnica de distracci\'on:} Alteraci\'on de cifras\\
\textbf{is\_duplicate:} false\\
\textbf{retry\_count:} 0
}
\end{tcolorbox}

\subsubsection{Ejemplo Few-Shot del Prompt}

Los ejemplos few-shot (ejemplos reales incluidos en las instrucciones al modelo
para que aprenda el formato esperado) est\'an basados en ex\'amenes oficiales:

\begin{tcolorbox}[colback=agentB!3, colframe=agentB,
  title={\small Ejemplo few-shot --- Carta de Naciones Unidas},
  fonttitle=\bfseries\small, boxrule=0.5pt, arc=3pt]
{\small
\textbf{Pregunta:} De los siguientes, \textquestiondown cu\'al NO es uno de los Consejos
que se establecen como \'organos principales de las Naciones Unidas?\\[4pt]
\begin{tabular}{cl}
\textbf{A)} & El Consejo de Administraci\'on Fiduciaria \\
\textbf{B)} & El Consejo Econ\'omico y Social \\
\textbf{C)} & El Consejo Internacional de Justicia \quad \checkmark \\
\textbf{D)} & El Consejo de Seguridad \\
\end{tabular}\\[4pt]
\textbf{T\'ecnica:} Atribuci\'on incorrecta de funciones\\
\textbf{Tip:} El \'organo judicial principal de la ONU se denomina
\textit{Corte} Internacional de Justicia, no \textit{Consejo}. Las dem\'as
opciones s\'i son consejos definidos como \'organos principales en la Carta.
}
\end{tcolorbox}

\begin{tcolorbox}[colback=agentB!3, colframe=agentB,
  title={\small Ejemplo few-shot --- Art. 25 DUDH},
  fonttitle=\bfseries\small, boxrule=0.5pt, arc=3pt]
{\small
\textbf{Pregunta:} Se\~nale la proposici\'on INCORRECTA.\\[4pt]
\begin{tabular}{cl}
\textbf{A)} & Toda persona tiene derecho a un recurso \\
    & efectivo ante los tribunales nacionales \\
\textbf{B)} & Toda persona tiene derecho a seguros en caso \\
    & de desempleo por \textit{cualquier} circunstancia \quad \checkmark \\
\textbf{C)} & Toda persona tiene derecho de acceso a las \\
    & funciones p\'ublicas de su pa\'is \\
\textbf{D)} & Toda persona tiene deberes respecto a la \\
    & comunidad \\
\end{tabular}\\[4pt]
\textbf{T\'ecnica:} Eliminaci\'on de excepciones legales\\
\textbf{Tip:} El Art. 25 especifica que el derecho a seguros se activa solo
por circunstancias \textit{independientes de la voluntad}. La opci\'on B elimina
esta excepci\'on diciendo ``cualquier circunstancia'', lo que la hace incorrecta.
}
\end{tcolorbox}

\end{document}